\chapter{Structure Polynomials and the Horizon Derivative}
\label{ch:strukturpolynome}

% Register: D3, D4, D7, D10, S2, S3, S4, S7, S39, S40, S41, S42, S43.
% Sources: Spec. 2, 3, 4, 9, Appendices H and I.

\section{The shape acquires coordinates}

So far the ``shape'' of a Hori number was nothing but its digit
picture. Now we give it an algebraic form -- and an intuition
becomes a computational tool.

The key are the \emph{falling factorials}
\[
\ff X0=1,\qquad \ff Xd=X(X-1)\cdots(X-d+1),
\]
the polynomials that count ``ordered selection without
replacement''. Reading the digits of a Hori number from right to
left as coefficients in this basis produces its \emph{structure
polynomial}:
\[
a=(a_1,\ldots,a_n)
\quad\longmapsto\quad
P_a(X)=\sum_{d=0}^{n-1}a_{n-d}\,\ff Xd .
\]

\begin{satz}[Evaluation theorem; Register S2]
$V_n(a)=P_a(n+1)$.
\end{satz}

\begin{proof}
For $d=n-i$ we have $\ff{(n+1)}d=(n+1)!/(i+1)!$ -- exactly the
weight of place $i$. Substituting yields the value formula.
\end{proof}

A small computation with a large consequence: the value of a Hori
number is \emph{the evaluation of its shape at the place of its
horizon}. The seven in $H_4$, that is $0012$, has the shape
$P(X)=X+2$; its value is $P(5)=7$. Prepend a zero and the shape
remains $X+2$ -- only the evaluation point moves: $P(6)=8$. The
mysterious value shift caused by leading zeros is no longer a
riddle but simply this: \emph{the same polynomial, evaluated one
place further on.}

Two precautions. First: only nonnegative integer coefficients are
admissible, and in horizon $n$ only polynomials with $d+c_d\le n$
at all occupied places ($c_d>0$) --
that is the digit bound in polynomial language; the quantity
$\sigma(P)=\max_{c_d>0}(d+c_d)$ -- the maximum runs over the
occupied places, and for the zero polynomial we agree that
$\sigma(0)=0$ -- is called the structure horizon (Register D4).
Second: the polynomial alone is \emph{not} the whole Hori number.
$12$, $012$ and $0012$ all have the shape $X+2$, yet they are three
different elements -- only the pair of horizon and polynomial
carries everything.

\section{The horizon derivative}

How strongly does a leading zero change the value? The answer is
the discrete derivative, known in analysis as the forward
difference:
\[
\Delta P(X)=P(X+1)-P(X),
\qquad
\Delta\,\ff Xd=d\,\ff X{d-1}.
\]
The second formula is the discrete counterpart of
$(x^d)'=d\,x^{d-1}$ -- the falling factorials are to the calculus
of differences what the powers are to differentiation.

\begin{satz}[Register S7]
The value change under a structure-preserving extension is
$\Delta P(n+1)$. In particular, a leading zero is value-neutral
exactly when $\Delta P=0$, that is, when $P$ is constant -- these
are precisely the numbers at or above their eigenhorizon.
\end{satz}

\begin{beispiel}
The shape $X+2$ (digits $\ldots12$) has derivative $1$: its values
climb by exactly one with each move, $5,6,7,8,\ldots$ The shape
$\ff X2$ -- that is the twelve in $H_3$, digits $100$ -- has
derivative $2X$: its values jump $12, 20, 30, 42,\ldots$, with
increments $8, 10, 12$ that grow themselves. A constant such as
the shape $7$ has derivative $0$ and never stirs.
\end{beispiel}

With this, the eigenhorizon from Chapter~\ref{ch:erweiterungen} is
no longer an intuitive notion but an equation: the \emph{shapes}
with horizon derivative zero are exactly the constant polynomials
-- read as shapes, the natural numbers. (Strictly speaking: as
elements, the constants form many horizontal representatives of
the same value; only in the shape-preserving limit space of
Chapter~\ref{ch:erweiterungen} does this become exactly one copy
of $\Nn$.) Everything else in the Hori space has a genuine
``velocity'': a shape of degree 1 grows by a constant amount per
horizon step, degree 2 accelerates, and so on. The degree of the
polynomial is the \emph{Hori dimension} of a shape (Register D7).

\section{Length times length is area -- now exact}

Dimension behaves additively under multiplication:
$\dim(PQ)=\dim P+\dim Q$ for $P,Q\ne0$. What is a mnemonic at school --
length times length gives area -- is here a theorem about degrees.
And one can compute with shapes directly. In $H_3$ the four has
the shape $X$ (digits $010$), the five the shape $X+1$ (digits
$011$). Their shape product is
\[
X(X+1)=\ff X2+2X
\]
(since $X^2=\ff X2+X$), hence the digit string $120$ -- and
indeed: $120$ in $H_3$ has the value $1\cdot12+2\cdot4+0=20$.
The multiplication $4\cdot5=20$ has happened here without a single
carry, purely by combining the shapes. The product formula of the
falling factorials carries a pretty side meaning along the way:
$X\cdot X=\ff X2+X$ splits all ordered pairs into those with
distinct and those with equal entries.

\section{What operations cost: the horizon calculus}
\label{sec:horizontkalkuel}

Every shape needs a horizon large enough for it -- the quantity
that measures this is the structure horizon $\sigma$. With it,
every operation acquires a \emph{price}: how much capacity must
the target space have so that the result is guaranteed to fit
there? For an operator $T$ we ask for the worst case
\[
\Gamma_T(r,s)=\max\{\sigma(T(P,Q)) \mid \sigma(P)\le r,\ \sigma(Q)\le s\}
\]
(Register D10; for unary operators, correspondingly
$\Gamma_T(r)$). At first sight a monster: already for $r=s=8$ more
than a hundred billion pairs would have to be checked. One
suffices.

\begin{satz}[Extremal principle; Register S39]
In every horizon $r$ there is the \emph{maximally filled} shape
$M_r$ -- every digit at the limit,
$M_r(X)=\sum_{d<r}(r{-}d)\,\ff Xd$; we also call it the
\emph{maximal polynomial} of the horizon.
For every operator that answers growing coefficients with growing
coefficients (all those composed from addition, multiplication
and $\Delta$ do), we have
\[
\Gamma_T(r,s)=\sigma\bigl(T(M_r,M_s)\bigr).
\]
\end{satz}

The worst case among untold billions of candidates is always the
same: all digits full. The proof (Appendix~\ref{ch:anhang-beweise})
is one line of monotonicity -- but the consequence is a
\emph{calculus}: addition costs exactly $r+s$; multiplication costs
$\beta(r,s)$, the product height, factorially growing at its core,
which we shall meet again in Chapter~\ref{ch:zaehlungen} as a
counting sequence. And the derivative delivers the biggest
surprise of the calculus:

\begin{satz}[Difference overflow; Register S40]
For $r\ge2$,
\[
\max_{\sigma(P)\le r}\sigma(\Delta P)
=\Bigl\lfloor\tfrac{(r+1)^2}{4}\Bigr\rfloor-1 .
\]
\end{satz}

Pause for a moment: $\Delta$ \emph{lowers} the degree -- and still
lets the capacity swell quadratically. Already at $r=8$ a single
difference lifts the structure horizon from 8 to 19. The witness
is a shape with a single digit in the middle: the derivative
multiplies each coefficient by its degree, and the most expensive
compromise between ``high degree'' and ``large coefficient'' lies
exactly in the middle -- hence the term
$\lfloor(r{+}1)^2/4\rfloor$, the largest product of two numbers
with a fixed sum. The lesson: the degree measures the
\emph{dimension} of a shape, $\sigma$ measures its
\emph{capacity}, and the two can jump in opposite directions. An
ordinary degree analysis is blind to this effect.

One operation is still missing from the calculus, the most
delicate one: \emph{substituting} one shape into another,
$P\circ Q$. That this is allowed at all is itself a theorem -- it
is by no means clear from the outset that substitution again
produces nothing but admissible, integral, nonnegative digits.

\begin{satz}[Composition closure; Register S41]
Substituting a shape into a shape yields a shape again.
Substitution, too, obeys the extremal principle:
$\Gamma_\circ(r,s)=\sigma(M_r\circ M_s)$.
\end{satz}

The proof (Appendix~\ref{ch:anhang-beweise}) is the combinatorial
centrepiece of this chapter: one reads $Q$ as a stock of
blueprints -- each digit $a_k$ provides $a_k$ templates, to be
mounted on $k$ distinct slots --, then $P\circ Q$ counts tuples
of such assemblies, and the required divisibility by $m!$ comes
from the fact that relabelling the slots cannot fix any choice of
assemblies. The cost of substitution, however, is of an entirely
different calibre from everything so far: the diagonal
$\Gamma_\circ(r,r)$ begins with
\[
1,\quad 4,\quad 23,\quad 6541,\quad 477\,149\,751,\quad\ldots
\]
-- substitution multiplies the degrees \emph{and} makes the
coefficients explode; this is the ninth and by far the
fastest-growing counting sequence of this project
(Chapter~\ref{ch:zaehlungen}). The most harmless special case,
the mere shift $P(X)\mapsto P(X{+}1)$, costs exactly
$r+\lfloor r^2/4\rfloor$ -- to the point exactly as much as one
derivative one horizon higher (Register S42), a small duality the
calculus throws in for free.

\section{What the shapes count: the counting width}
\label{sec:zaehlbreite}

To close the chapter, the question one should always put to an
algebraic object: does it count something? The answer arrived as
an external classification of the project and is so far
Horimetrik's best bridge into an established field --
\emph{counting complexity}.

The key is the basic meaning of the falling factorial:
$n^{\underline d}$ counts the ordered $d$-tuples of
\emph{distinct} objects from a stock of $n$. If now $n$ is the
number of solutions of some search problem -- in computer science
one says: of its \emph{witnesses} --, then a shape
$P=\sum_d c_d\ff Xd$ reads as a blueprint for new counted objects:
$c_0$ fixed special objects, $c_1$ colours for each individual
solution, $c_2$ variants for each ordered pair of distinct
solutions, and so on. $P(n)$ is exactly the number of these
coloured witness tuples. This harmless observation has a precise
consequence:

\begin{satz}[Hori \#P theorem; Register S43]
If $f$ is the witness-counting function of an arbitrary NP
problem and $P$ a shape, then $P\circ f$ is again the
witness-counting function of an NP problem. Every shape is thus a
\emph{closure operator} of the counting class \#P.
\end{satz}

The proof (Appendix~\ref{ch:anhang-beweise}) is exactly the
blueprint reading: guess the degree, guess the colour, guess a
tuple of distinct witnesses. Honesty first: the \emph{qualitative}
side of this statement is known -- polynomials with nonnegative
integer coefficients in the \emph{binomial} basis are called
\emph{binomial-good} in the literature and are characterized as
the (relativizing) \#P closure operations. Because of
$\ff Xd=d!\binom Xd$ our shapes are a \emph{proper subclass} of
these: divisibility by $d!$ is exactly the difference between
ordered tuples and unordered selections. What the literature
apparently does \emph{not} have is the quantitative layer -- and
Horimetrik supplies it for free:

Suddenly all the quantities of this chapter mean something. The
\emph{addition} of shapes is the disjoint union of two counting
constructions. \emph{Multiplication} couples two constructions,
and the coefficients of its product formula count exactly the
overlap patterns of two groups of witnesses -- so the growth of
$\beta$ measures how the variety of overlaps explodes under
coupling. The \emph{horizon derivative} counts which new
structures a single additional witness makes possible
($d$ positions times remaining tuples: exactly $d\,\ff n{d-1}$) --
and the difference overflow says: a new witness lowers the tuple
order, yet can inflate the variety of variants quadratically.
And $\sigma$ itself? Draw the coefficient board -- column $d$ with
$c_d$ boxes --, then $\sigma(P)=\max(d+c_d)$ is the smallest
\emph{staircase} into which the board fits: the maximum of
``how many witnesses at once'' and ``how many variants of them''.
We call this quantity the \emph{counting width} of a
construction. Every horizon $r$ is thus a finite, completely
enumerable family of exactly $(r{+}1)!$
\#P closure operators, and the horizon calculus of
Section~\ref{sec:horizontkalkuel} gives the optimal figures for
how the counting width grows under uniting, coupling, extending
and substituting.

Two sentences of demarcation, so that no false magic arises: in
this neighbourhood lives a weighty open conjecture (even its
simplest version would imply P$\,\ne\,$NP), and our operators lie
expressly on the \emph{well-understood} side of this landscape --
Horimetrik solves nothing here. What it contributes is more
modest, and perhaps usable for precisely that reason: turning an
allowed-or-not-allowed, for the first time, into a \emph{how
wide}.

\begin{oberflaeche}
What this chapter has delivered, in one sentence: the shape of a
Hori number is a polynomial, its value is that polynomial
evaluated at the place of its horizon, and the puzzling effect of
leading zeros is simply the derivative of this polynomial. The
ordinary numbers are the shapes with derivative zero -- the only
inhabitants of the Hori space who do not care where they live.
All the others move when their world grows. Every arithmetic
operation has a price in capacity that can be stated exactly, and
it is always the same extreme object that pays it: the shape
filled to the brim. And at the end it turned out that the shapes
\emph{count} something -- coloured groups of distinct solutions
of a search problem -- and that $\sigma$ measures the width of
these counting constructions. In the next chapter we let these
moving objects do arithmetic -- and run into the question of what
should then happen to their past.
\end{oberflaeche}
